\subsection{Where metabolic memory lives}\label{sec:memory}

Two established results, put side by side, force a structural
conclusion. First, $66\%$ of closed genotype cycles retain
flux-state drift after full restoration (\S\ref{sec:computational}):
the flux state after an $A \to AB \to B \to A$ cycle is not the
initial state, with a median residual drift of $\sim 110$ $\ell_1$
units. Second, the protein layer does not carry the association:
protein fold change explains at most $0.7\%$ of the curvature
variance ($r = -0.083$, $R^2 = 0.007$, $n = 366$;
\S\ref{sec:e27}), while transcripts on the same genes carry $r =
+0.420$; under the locked $\kmu$ metric the picture is unchanged
(protein $r = -0.098$, $p = 0.06$; transcript $r = +0.339$). A
state difference of median size $110$ cannot be stored
in a layer that moves by less than one percent of the signal:
enzyme abundance is quantitatively excluded as the carrier of
metabolic hysteresis. Transcripts carry the association but are
instructions, not state --- the menu, not the meal. The persistent
state variable must therefore be fast and post-translational:
metabolite pools and their modification states.

This deduction tells the wet-lab investigator where not to look
(enzyme titers) and where to look. \citet{kochanowski2013}
document metabolite-level passive coordination on precisely this
class of timescales, and the wall coordinates of
\S\ref{sec:walls} name the candidate pools: the
F6P/GAP/PEP forks that carry most of the curvature mass. A
testable prediction follows. In public metabolomics time series
across nutrient transitions, branch-metabolite pools should show
slope kinks at wall crossings while dedicated-metabolite pools
respond smoothly; and the same pools, not dedicated ones, should
carry the non-reverting offset after a restoration cycle. The
branch-metabolite table is thus the shared substrate of two of
this paper's claims --- the chemical coordinates of the walls and
the carrier of their memory.
